%----------
%	WORK METHODOLOGY AND PLANNING
%----------

Para estructurar el desarrollo de la investigación y garantizar el cumplimiento de los objetivos, el proyecto se ha dividido en seis Paquetes de Trabajo (PT). A continuación, se detallan las tareas principales y los entregables asociados a cada uno de ellos:


\begin{itemize}
    \item \textbf{PT0. Coordinación, comunicación y diseminación}
    \begin{itemize}
        \item Reuniones periódicas con el equipo de IRIS\_IAMEDIA.
        \item Supervisión académica y de RTVE.
        \item Documentación y difusión de avances.
        \item Preparación de publicaciones para congresos, conferencias y revistas científicas de ámbito nacional e internacional.
        \item Participación en eventos de divulgación y jornadas de comunicación.
        \item \textbf{Entregable:} Entrega final y defensa del TFM en RTVE y en UC3M.
    \end{itemize}
    
    \item \textbf{PT1. Diseño de la arquitectura agéntica}
    \begin{itemize}
        \item Definición de las agentes especializadas y sus roles.
        \item Diseño del bucle agéntico de razonamiento y acción, con verificación de evidencias.
        \item Especificación del formato de salida estructurado.
        \item Esquemas de interacción entre agentes, herramientas y niveles de contexto.
    \end{itemize}
    
    \item \textbf{PT2. Diseño y evaluación de Agent Skills y prompts}
    \begin{itemize}
        \item Formalización de la metodología experta como \textit{skills} cargables.
        \item Creación de plantillas y estrategias de \textit{prompting}.
        \item Experimentos y pruebas de consistencia, con ajustes iterativos.
    \end{itemize}
    
    \item \textbf{PT3. Representaciones semánticas y retrieval}
    \begin{itemize}
        \item Implementación de RAG.
        \item Índices de similitud.
        \item Evaluación de precisión contextual.
    \end{itemize}
    
    \item \textbf{PT4. Evaluación comparada de modelos}
    \begin{itemize}
        \item Pruebas con modelos comerciales por API y modelos locales de pesos abiertos.
        \item Métricas de rendimiento y de acuerdo.
        \item Identificación de fortalezas y limitaciones.
    \end{itemize}

    \newpage 
    
    \item \textbf{PT5. Prueba de concepto, conclusiones y transferencia}
    \begin{itemize}
        \item Integración con el proyecto IRIS\_IAMEDIA.
        \item Directrices de integración editorial.
        \item Conclusiones y recomendaciones finales.
    \end{itemize}
\end{itemize}



\begin{landscape}
\begin{figure}[H]
	\ffigbox[\FBwidth]
	{\caption[Cronograma de trabajo (diagrama de Gantt)]
 {Cronograma de trabajo (diagrama de Gantt)}}
 % [Aquí ponemos como queremos que aparezca en el índice, sin la referencia]{Aquí como queremos que aparezca en el documento, con la referencia}
	{\includegraphics[scale=0.6]{Plantilla_TFG_ingles_2019/imagenes/cronograma_iris_blanco.pdf}}
\end{figure}
\end{landscape}

\newpage





